\chapter{课题研究目标、研究内容和拟解决的关键性问题}

\section{研究目标}

为填补学术界在多FPGA硬件仿真平台方面的研究空缺，本课题提出REMUS（Scalable Multi-FPGA Replayable Emulator System），一个面向多层次（片内、跨片、跨服务器）与多粒度（模块级与网表级）的多FPGA硬件仿真的开源解决方案。REMUS面向三个核心目标：可调试性、通用性、自动化。为实现这些目标，本课题在多FPGA系统部署的三大步骤——电路划分、结构集成、互联体系——中引入系统性创新，包括基于扫描链的硬件仿真技术路线、结合延迟不敏感与分时复用优势的LI-TDM跨边界接口、面向通用RTL的全自动划分与集成流程，以及基于AXI-Stream协议的流式互联通信抽象。

REMUS首先是一个周期精确、比特精确的多FPGA硬件仿真平台，能够在软件仿真环境中精确重放硬件任意周期的内部状态。其次，它提供了从RTL输入到可综合Verilog输出的全自动化流程，仅需指定划分数量即可完成完整的划分与集成，无需用户介入。再次，REMUS是一个灵活的多粒度解决方案：设计既可在模块级或网表级进行划分，也可基于资源评估动态选择合适粒度；同时支持对部分模块选择性插入扫描链以降低资源开销。最后，REMUS是跨层次通用的互联体系，支持三种部署方式：可重构区域的片内互联、点对点跨FPGA互联，以及跨服务器网络互联。

\section{关键性问题}

\begin{itemize}

    \item 电路划分：自适应粒度的全自动划分流程。实现支持自适应粒度的自动化电路划分框架，需要集成超图划分算法，并保持高度可扩展性，使设计在模块级与网表级均可进行有效分割。
          
    \item 结构集成：结合扫描链与 LI-TDM的统一集成技术。基于LI-TDM的轻量级跨边界接口与扫描链结构，对划分后的子电路进行统一集成，从而实现对内支持任意同步电路的精准划分，对外支持多种多FPGA互联架构，并为硬件仿真提供可重放的扫描链数据路径。
          
    \item 互联体系：统一可扩展的跨FPGA通信和主机通信。构建基于上述跨边界接口的多FPGA流式互联通信抽象，与物理信道解耦，可部署于片内互联、跨片点对点链路以及跨服务器网络等多层次的系统中。同时在有限的接口下构建可扩展的主机通信链路，实现一对多的控制信号和检查点数据的传输。
          
    \item 系统验证：在自建多FPGA平台上的大规模设计验证。在自主搭建的多FPGA平台上部署多种处理器或SoC的验证场景，验证所提出方法在运行性能、可扩展性与可调试性，并给出系统级评估结果。
          
\end{itemize}

\section{研究内容}

本节介绍REMUS的系统框架与设计思路。REMUS面向跨节点的多FPGA部署，旨在克服现有方案在可调试性、扩展性与自动化方面的不足。其整体流程遵循上述的多FPGA系统部署的步骤，即：电路划分、结构集成、互联体系三个核心阶段。用户只需将自定义硬件逻辑接入REMUS提供的标准化顶层接口（包括AXI存储器、串口、时钟等基础外设），无需额外修改即可进入完全自动化的处理流程。

REMUS的运行机制如图\ref{fig::remus-framework}所示：首先对输入设计执行自动化的电路划分；随后对每个划分子模块进行结构集成，包括跨FPGA数据接口与扫描链集成；最后，在多级物理架构上构建统一的通信体系，并通过运行时软件实现协同执行、暂停恢复与波形重放等能力。

\begin{figure}[htb]
    \centering
    \includegraphics[width=\columnwidth]{figure/REMUS-framework.png}
    \bicaption{REMUS框架示意图}{REMUS framework}
    \label{fig::remus-framework}
\end{figure}

\subsection{电路划分：自适应粒度的全自动划分流程}

REMUS 的第一阶段是自动化电路划分。系统基于综合工具解析并初步综合用户提供的Verilog设计，在保留层次结构的同时将内部逻辑转化为中间层原语。随后，REMUS执行划分预处理，包括模块资源评估、跨模块连线统计以及对资源占用较大的模块进行递归展开。这一预处理步骤是实现自适应粒度划分的关键，使系统能够在模块级与网表级间灵活切换，并依据资源配置自动选取合适粒度。

在此基础上，框架将预处理后的设计结构自动转换为标准化的超图表示，调用超图划分算法获得划分映射。系统在该阶段预留接口，可便捷地集成其他划分算法。最终，REMUS将划分结果写回至Verilog格式，为每个子模块插入划分属性并生成结构化的划分输出文件。划分阶段的输入与输出均保持Verilog形式，以便独立集成到其他多FPGA流程之中。

\subsection{结构集成：基于LI-TDM与扫描链的统一集成技术}

划分完成后，REMUS自动进入结构集成阶段。该阶段包含划分集成与逻辑集成两部分。

在划分集成部分，框架首先对划分产生的互联边界进行一次统一跨边界接口，将线网连接转换为协议化的数据传输接口。REMUS采用LI-TDM跨边界接口方案，它结合了延迟不敏感（Latency-Insensitive， LI）通信与分时复用（Time-Division Multiplexing， TDM）的优势，对外提供延迟不敏感接口，从而在跨FPGA通信中显著降低时序约束；对内以TDM机制驱动状态机，天然支持任意同步时序电路的划分，即便在组合逻辑存在依赖的划分边界也能避免死锁。

在逻辑集成部分，REMUS对内部所有时序器件进行统一的时钟门控改写、时钟域分析与扫描链插入，生成可精确重构电路状态的扫描链端口。然而原始扫描链控制器与LI-TDM控制器语义冲突，无法直接复用。为此，REMUS设计了新的统一控制逻辑，实现多FPGA的同步启停、扫描链双向数据传输以及周期级精确状态捕获。整个流程完全自动化，由综合工具解析划分结果并生成最终可综合的Verilog，无需用户干预。


\subsection{互联体系：统一可扩展的跨FPGA通信和主机通信}

通信阶段涉及FPGA间的运行时数据传输以及主机与FPGA间的控制通信。由于划分将电路逻辑拆解至不同FPGA，运行时中间结果无法直接送达目标寄存器，必须依赖跨FPGA通信路径进行传输。为实现通用性，REMUS为所有划分子设计提供统一的数据互联接口，使其与具体物理信道解耦。

基于多FPGA系统中对可靠数据包与动态路由的需求，REMUS采用AXI-Stream协议构建抽象化的流式互联接口，并在片内互联、跨片链路以及跨服务器网络三种架构中实现同样的通信语义，展示其多层次通用性。

另一方面，作为周期精确的硬件仿真系统，REMUS还需支持软件侧对FPGA的统一启停控制、检查点传输与波形重放。单片系统REMU采用PCIe + DMA的直接主机通信方式，但扩展至多FPGA后面临主机接口数量受限与多片消息同步的问题。为此，REMUS引入分层控制路径，通过复用FPGA间现有物理链路减少主机直接通信需求，并依靠分布式集成逻辑实现多片间同步推进。

\subsection{方法优势与对比分析}

表~\ref{tab:compare}所示，REMUS在可调试性、通用性与自动化程度上，与现有学术系统相比，均有显著提升。其自动化划分流程支持可变粒度甚至自适应粒度选择，而FireAxe在划分阶段仍需手动指导，Hetero-ViTAL仅支持模块级粗粒度划分，DUOHUAFEN仅支持网表级划分。在集成层面，REMUS是首个将扫描链与划分接口统一化的方案，实现统一的启停控制；其划分集成基于LI-TDM，既支持任意同步时序电路，又向外提供延迟不敏感接口，适用于同步与异步多FPGA系统。相比之下，FireAxe的LI-BDN接口控制逻辑复杂，而Hetero-ViTAL的LI通道限制较强，无法应用于任意电路的划分。在通信体系上，REMUS提供物理信道无关的流式互联协议，可适应异构链路结构；而Hetero-ViTAL依赖特定云平台，DUOHUAFEN与FireAxe均基于同构通信结构，通用性有限。

\begin{table}[htb]
    \centering
    \caption{REMUS与相关工作的功能对比}
    \label{tab:compare}
    \begin{tabular}{*{6}{c}}
        \toprule
        框架名称  & FireAxe & Hetero-ViTAL & DUOHUAFEN & REMU   & REMUS  \\
        \midrule
        划分粒度  & 网表      & 模块           & 网表        & -      & 网表或模块  \\
        自动划分  & 否       & 是            & 是         & -      & 是      \\
        跨边界接口 & LI-BDN  & LI           & LI-TDM    & -      & LI-TDM \\
        硬件仿真  & 否       & 否            & 否         & 是      & 是      \\
        通信体系  & 同构通信    & 网络           & 同构通信      & PCIe直连 & 异构通信   \\
        \bottomrule
    \end{tabular}
\end{table}